
\chapter{Introdução}
\label{cap:intro}

Este relatório documenta a atividade da disciplina Tópicos Avançados em
Inteligência Artificial I, reproduzida a seguir na íntegra:

\begin{citacao}
\textbf{Ensamble \textit{(sic)}: clusterização de clientes com Regras de Associação}

Combinar clusterização de clientes com Regras de Associação é uma estratégia
extremamente poderosa para análises de CRM e marketing. A ideia central é
simples: em vez de aplicar as regras de associação em toda a base de
transações (o que pode gerar milhares de regras genéricas e pouco
acionáveis), você segmenta os clientes primeiro e, em seguida, descobre
padrões de compra específicos para cada grupo.

Portanto, de posse do DW Organizacional ou Datamart construído para esse fim,
povoar o BD e aplicar o conceito de ensemble (aprendizagem por conjuntos):

1) Utilizar dados para clusterizar ou classificar os clientes

2) De posse das compras por grupo/classe, aplicar as Regras de Descoberta de
Associação entre produtos.

Obs: Para o trabalho pode-se utilizar as ferramentas das bibliotecas Python.

Entregáveis: Código e relatório técnico do processo.
\end{citacao}

O enunciado chama de \textit{ensemble} a combinação, em sequência, de duas
técnicas não supervisionadas: segmentar os clientes antes de minerar regras
de associação, em vez de minerar sobre toda a base de uma só vez. É um uso
mais amplo do termo, que na literatura designa em geral a combinação das
predições de vários modelos \cite{dietterich2000}.

\section{Ambiente construído e adaptação do domínio}
\label{sec:ambiente-dominio}

Para esta atividade foi construído, do zero, um ambiente sobre os dados
abertos de sinistros de trânsito da \sigla{PRF}{Polícia Rodoviária Federal}
\cite{prfdatatran}: \sigla{DW}{Data Warehouse} Organizacional (Inmon
\cite{inmon2005}), Data Mart em estrela (Kimball \cite{kimball2013}),
\sigla{ETL}{Extract, Transform, Load} via Airflow e \sigla{BI}{Business
Intelligence} via Metabase (\autoref{cap:ambiente}). Sobre o Data Mart
povoado foi aplicado o \textit{ensemble} (\autoref{cap:metodologia}), com
resultados no \autoref{cap:resultados}.

O domínio de origem (acidentes de trânsito) não tem o conceito de
\textit{cliente} comprando \textit{produtos}. Os dois conceitos foram
adaptados ao domínio disponível, preservando a mesma estrutura do problema
(segmentar, depois minerar regras por segmento):

\begin{quadro}[!ht]
  \centering
  \small
  \caption{Adaptação dos conceitos do enunciado ao domínio PRF-DATATRAN}
  \label{qua:adaptacao}
  \begin{tabular}{|p{2.6cm}|p{5.2cm}|p{6.2cm}|}
    \hline
    \textbf{Enunciado} & \textbf{Adaptado para} & \textbf{Justificativa} \\
    \hline
    Cliente &
    Município (\texttt{nome\_municipio} + \texttt{sigla\_uf}) &
    Unidade geográfica que concentra ocorrências ao longo do tempo; acumula
    um \textit{perfil de risco} (letalidade, gravidade, volume, condições
    predominantes), como um cliente acumula um perfil de compras \\
    \hline
    Produto (item da cesta de compra) &
    Cada classificação de uma ocorrência (causa, tipo de acidente, condição
    meteorológica, fase do dia, tipo de pista, traçado da via, uso do solo,
    sentido da via, classificação quanto a vítimas) &
    Uma ocorrência é tratada como uma \textit{cesta} com até 9 itens, um por
    tipo de classificação \\
    \hline
  \end{tabular}
  \fontequa{Elaborado pelo autor.}
\end{quadro}

Todo o trabalho (código-fonte e este relatório) está no repositório
\url{https://github.com/coutinhonobre/dw_trabalho_apresentacao}
\cite{repotrabalho}. O relatório usa o modelo \LaTeX\ oficial do IFG,
\textsf{classe-ifg}, do Prof. Dr. Raphael de Aquino Gomes
\cite{classeifg}.
